\section{Introdução}
Uma gargalhada pode aparecer na escrita como \forma{ah! ah! ah!}, \forma{quiá-quiá}, \forma{cá-cá}, \forma{quá-quá}, \forma{hahaha}, \forma{rsrsrs} ou \forma{kkkk}. Na comunicação digital, o repertório também inclui emoticons, emojis, sequências heterogêneas de letras, imagens e respostas verbais como \forma{rindo}. A convergência funcional dessas formas é evidente: todas podem participar da representação ou da negociação do humor. Sua relação histórica exige outra espécie de evidência. Compartilhar um referente ou uma motivação sonora não demonstra que uma forma descenda de outra.

A pergunta central deste artigo é: \emph{quais evidências permitem distinguir continuidade histórica, convergência icônica e inovação gráfica na representação do riso em português?} A questão ganha relevância quando explicações de circulação pública ligam a escrita digital brasileira a passagens de José de Alencar, Machado de Assis ou Eça de Queirós, sem mostrar os elos intermediários. Um registro oitocentista pode ser um antecedente documental da representação escrita do riso, mas isso não o converte automaticamente em origem etimológica de uma convenção digital.

O problema inverso é igualmente importante. Seria incorreto rejeitar essa genealogia afirmando que onomatopeias, por definição, não possuem história, não podem ser transmitidas ou nunca se integram à gramática. O caráter imitativo de uma expressão não a impede de ser aprendido, convencionalizado, emprestado, lexicalizado ou transformado. A literatura recente mostra a participação da experiência linguística na percepção de iconicidade \cite{iida2024}. A alternativa entre ``evolução real'' e ``mera onomatopeia'' precisa ser substituída por dimensões que podem coexistir.

O material que motivou o estudo reúne explicações conversacionais sobre mudança fonética, semântica e gramatical; exemplos como \forma{olho}, \forma{você} e \forma{vilão}; comparações entre idiomas; alegações sobre a internet brasileira dos anos 1990 e 2000; e panoramas decenais de livros, jornais e quadrinhos. Inclui frases apresentadas como citações literais, às vezes acompanhadas de datas precisas, mas sem edição, página ou reprodução documental. Essa combinação torna necessária uma crítica de fontes: a precisão aparente de uma data não substitui a recuperação do documento.

O objetivo geral é construir uma interpretação linguística e documental das representações do riso, preservando as perguntas do material de partida e examinando criticamente suas respostas. Os objetivos específicos são: definir mudança, derivação, parentesco, iconicidade e coincidência; revisar exemplos etimológicos; recuperar testemunhos com localizadores; auditar cronologias; fundamentar estatisticamente a ampliação do corpus; e disponibilizar instrumentos de reprodução. A repetição de uma afirmação em várias respostas não produz novas fontes independentes.

Emprega-se \emph{português brasileiro} para a variedade investigada, sem pressupor que uma única gramática normativa explique todos os usos digitais. ``Gramática da língua brasileira'', no material inicial, é interpretada como interesse pela organização do português em uso no Brasil. Distinguem-se \emph{sorriso}, expressão facial que pode ser silenciosa, \emph{riso}, categoria mais ampla, e \emph{gargalhada}, manifestação vocal intensa. O objeto é o riso escrito: um marcador digital de humor não comprova que o escrevente tenha rido fisicamente.

Os fundamentos de linguística e de fonética e fonologia apoiam-se em obras de referência \cite{martelotta2008,seara2015}. Estudos de 2024 e 2025 complementam os fundamentos em iconicidade e comunicação digital \cite{iida2024,badrakh2024,kempe2025}.

\section{Problema historiográfico e hipóteses concorrentes}
\subsection{Anterioridade, descendência e equivalência funcional}
Há quatro proposições diferentes: uma forma aparece antes de outra; duas formas representam riso; usuários de uma tiveram contato com a outra; uma convenção foi transformada a partir da outra. As duas primeiras são compatíveis com múltiplas histórias. A terceira exige evidência de circulação; a quarta, uma explicação de transmissão apoiada em documentação e contexto.

O enunciado ``o \forma{kkkk} nasceu em \textit{Til}, em 1872'' funde essas proposições. A presença de \forma{quiá} não equivale à presença do caractere repetido \forma{k}. Por outro lado, ``o \forma{quiá} morreu nos livros'' também excede a evidência sem pesquisa suficientemente ampla de cartas, periódicos, teatro, quadrinhos e práticas escolares. Obras impressas podem funcionar como canais de transmissão; uma tradição escrita não depende de passagem obrigatória pela fala cotidiana.

Consideram-se três hipóteses: H1, transmissão ou adaptação de convenções preexistentes; H2, recriação independente motivada por iconicidade e propriedades gráficas; H3, origem múltipla ou convergente, com transmissão local, reinvenção e difusão em redes. Nenhuma recebe exclusividade por definição. A semelhança entre \forma{quiá}, \forma{cá} e o nome da letra \forma{k} torna relações plausíveis, mas plausibilidade não demonstra causalidade.

\subsection{Primeira ocorrência encontrada e origem absoluta}
Uma ocorrência datada estabelece que a forma estava disponível naquele documento ou na interação nele reproduzida. Oferece, se o testemunho for confiável, um limite superior para a emergência: a inovação ocorreu até aquela data. Não informa necessariamente onde, por quem ou quanto antes. ``Primeira ocorrência encontrada neste levantamento'' difere de ``primeira ocorrência da língua''.

A conservação de registros é desigual. Conversas privadas de ICQ ou MSN, canais de IRC, BBS e salas públicas nem sempre foram arquivados. Um log preservado exige validação de data, integridade e procedência. A ausência de \forma{kkkk} em um conjunto de conversas entre 1995 e 1997, caso recuperado, seria resultado sobre aquele conjunto, não prova de inexistência no Brasil. Sem corpus identificado, a própria alegação de ausência permanece sem base verificável.
